% Chapter 5: Webbase Ontology Language (WBOL)

\chapter{Webbase Ontology Language (WBOL)}
\label{chap:wbol}

\begin{quote}
\textit{"The limits of my language mean the limits of my world."} \\
— Ludwig Wittgenstein
\end{quote}

\wbol{} is formally defined using a standard \textbf{XML Schema (XSD)}. A webbase is an XML document validated against this schema; the schema's \texttt{key}/\texttt{keyref} constraints enforce referential integrity across the GUIDs that tie the document together (homepages, mainpages, shortcut targets, option targets, and role references), something a plain data-interchange format cannot guarantee on its own.

\section{WBOL Ecosystem}
\label{sec:wbol-ecosystem}

While \wbol{} serves as the declarative language for describing portals, it operates within a comprehensive ecosystem of associated technologies and APIs that enhance its capabilities:

\subsection{Layout API}
\label{sec:layout-api}

\wbol{} includes an associated Layout API that provides programmatic control over content presentation and styling. This API works in conjunction with the declarative \texttt{STWLayout} elements to enable:

\begin{itemize}
\item Dynamic layout adjustments based on content type and user preferences
\item Responsive design patterns that adapt to different screen sizes and devices
\item Integration with external CSS frameworks and design systems
\item Runtime modification of layout properties based on data characteristics
\end{itemize}

The Layout API ensures that content presentation remains consistent across different contexts while allowing for necessary flexibility in complex enterprise environments.

\subsection{Text Preprocessor}
\label{sec:text-preprocessor}

The \wbol{} text preprocessor is a powerful engine that processes textual content before rendering, enabling:

\begin{itemize}
\item Placeholder substitution using the \textbf{Webbase Placeholders Language (\wbpl{})}
\item Localization and internationalization support
\item Dynamic text generation based on context and user data
\item Integration with external content management systems
\item Markdown and other markup language processing
\end{itemize}

This preprocessor is particularly crucial for the \texttt{command} attributes in \texttt{STWContent} elements, where it replaces placeholders with actual values from various sources before command execution.

\subsection{Visual Component Library (VCL)}
\label{sec:vcl}

The Visual Component Library provides a standardized set of UI components that correspond to the different content categories and subtypes. The VCL ensures:

\begin{itemize}
\item Consistent visual presentation across different portal implementations
\item Accessibility compliance through standardized component behavior
\item Cross-platform compatibility for various front-end frameworks
\item Extensibility for custom component development
\end{itemize}

\subsection{Content Management Integration}
\label{sec:cms-integration}

\wbol{} is designed as a fundamental language for Content Management Systems (CMS), providing:

\begin{itemize}
\item Structured content definition that separates presentation from data
\item Version control capabilities for portal configurations
\item Workflow management for content approval processes
\item Integration points for external content repositories
\end{itemize}

This integration capability makes \wbol{} particularly suitable for enterprise environments where content management workflows are critical to business operations.

\section{Spin the Web Studio}
\label{sec:studio}

The third and final component of Spin the Web is the \textbf{Spin the Web Studio}. Alongside the \textbf{\wbol{}} and the \textbf{Web Spinner}, it completes the framework. The Spin the Web Studio is a specialized \webbaselet{} engineered for editing \webbase{s}. To use it, you simply add the \webbaselet{} to the \webbase{} you wish to edit, enabling direct, in-place modification.

\section{The Common Element Structure}
\label{sec:stwelement-base}

Every \wbol{} element---\texttt{webbase}, \texttt{area}/\texttt{webbaselet}, \texttt{page}, and \texttt{content}/\texttt{shortcut}---shares a common shape: a \texttt{guid} identity attribute, a localized \texttt{name}, optional \texttt{visibilities} and \texttt{documentation}, and, where applicable, a \texttt{children} element. In the runtime (see \cref{chap:web-spinner}), these XML elements are parsed into the \texttt{STWPortal}, \texttt{STWArea}, \texttt{STWPage}, and \texttt{STWContent} in-memory classes respectively; the class hierarchy is an implementation detail of the \webspinner{}, not a second schema. The full XSD is in Appendix~\ref{app:wbol-xml-schema} (see Appendix: WBOL XML Schema Reference).

\subsection{Property Description}

\begin{description}
\item[\textbf{guid}]: A unique identifier for the element, encoded as an attribute (GUID). Each element kind uses a distinct prefix (e.g., \texttt{w...} for webbases, \texttt{a...} for areas, \texttt{p...} for pages, \texttt{c...} for content), which the schema uses to type-check references.
\item[\textbf{element name}]: The specific kind of the element is given by its XML tag itself---\texttt{webbase}, \texttt{area}, \texttt{webbaselet}, \texttt{page}, \texttt{content}, or \texttt{shortcut}---rather than by a separate discriminator attribute.
\item[\textbf{name}]: A localizable name for the element: a \texttt{name} element containing one or more \texttt{text} children, each qualified with an \texttt{xml:lang} attribute (e.g., \texttt{en}, \texttt{it}) and, for \texttt{name}, a required \texttt{slug} attribute. Example: \lstinline[language=XML]|<name><text xml:lang="en" slug="rd">Research \& Development</text></name>|
\item[\textbf{slug}]: A localizable, URL-friendly identifier, carried as the \texttt{slug} attribute on each \texttt{name/text} entry (one slug per language). Slugs contain only lowercase letters, numbers, periods, and hyphens.
\item[\textbf{keywords}]: Localizable keywords for SEO, structured like \texttt{name} but without the \texttt{slug} attribute.
\item[\textbf{description}]: A localizable description of the element, structured like \texttt{keywords}.
\item[\textbf{visibilities}]: Defines the visibility rules for the element based on user roles: a \texttt{visibilities} element containing \texttt{role} entries, each pairing a role \texttt{id} with a \texttt{visibility} value (\texttt{true}, \texttt{false}, or a visibility function call). If a rule for a role is not defined, the visibility is determined by checking the parent element's visibility rules. This hierarchical check continues up to the root element. If no rule is found, the element is not visible by default.
\item[\textbf{children}]: A \texttt{children} element containing an ordered mix of nested \texttt{area}, \texttt{webbaselet}, \texttt{page}, \texttt{content}, and \texttt{shortcut} elements, as constrained by the parent's type.
\end{description}

\section{WBOL Element Types}
\label{sec:wbol-element-types}

This section describes the specialized element types available in \wbol{}.

\subsection{STWPortal}

\texttt{STWPortal} is a singleton element that represents the entire portal. It inherits from \texttt{STWElement} and acts as the root element of the portal structure.

\subsubsection{Property Description}

\begin{description}
\item[\textbf{langs}]: A list of supported languages for the portal. The first of the list is the default portal language.
\item[\textbf{datasources}]: Defines the data sources used by the portal.
\item[\textbf{mainpage}]: The GUID of the \texttt{STWPage} that serves as the main entry point for the portal.
\item[\textbf{version}]: A version string for the portal.
\end{description}

\subsubsection{Data Sources and Groups}
\label{sec:datasources-groups}

The \texttt{datasources} element within the webbase root provides a flexible framework for defining how the portal connects to and manages external data systems. Beyond simple connection definitions, \wbol{} supports advanced data organization concepts:

\begin{description}
\item[\textbf{Data Source Configuration}]: Each \texttt{datasource} element is configured with an \texttt{id}, a \texttt{queryLanguage}, and a \texttt{connection-string}. Supported data source types include:
\begin{itemize}
\item Relational databases (SQL-based)
\item NoSQL databases (document, key-value, graph)
\item REST APIs and web services
\item File-based data sources (CSV, JSON, XML)
\item Enterprise systems (ERP, CRM, HRM)
\end{itemize}

\item[\textbf{Data Groups}]: Data sources can be organized into logical groups that facilitate:
\begin{itemize}
\item Access control and security policies
\item Performance optimization through connection pooling
\item Load balancing across redundant data sources
\item Backup and failover configurations
\item Administrative grouping for different business units
\end{itemize}

\item[\textbf{Accessibility Integration}]: The data source framework includes provisions for accessibility compliance:
\begin{itemize}
\item Metadata enrichment for screen readers and assistive technologies
\item Alternative data representations for different accessibility needs
\item Command result formatting that supports accessibility standards
\item Integration with accessibility frameworks and tools
\end{itemize}
\end{description}

This comprehensive approach to data source management ensures that \wbol{} can integrate seamlessly with complex enterprise data environments while maintaining accessibility and security standards.

\subsection{STWArea}

\texttt{STWArea} represents a logical grouping of pages, analogous to a chapter in a book. It extends the base \texttt{STWElement}.

\subsubsection{Property Description}

\begin{description}
\item[\textbf{mainpage}]: The GUID of the \texttt{STWPage} that serves as the main entry point for this area.
\item[\textbf{version}]: A version string for the area.
\end{description}

\subsection{STWPage}

\texttt{STWPage} represents a single page in the portal. It extends the base \texttt{STWElement} but restricts its \texttt{children} to only \texttt{STWContent} elements.

\subsection{STWContent}

\texttt{STWContent} represents a piece of content on a page, parsed from a \texttt{content} (or \texttt{shortcut}) element. It extends the base structure but is not allowed to have any \texttt{children}. It adds several attributes for data binding and layout control.

\subsubsection{Property Description}

\begin{description}
\item[\textbf{type}]: The content's \texttt{type} attribute, fixed to one of the schema's content kinds (see Appendix~\ref{app:wbol-xml-schema}).
\item[\textbf{subtype}]: Specifies the type of content to be rendered, which determines the component used on the front-end. Possible values include \texttt{Text}, \texttt{Form}, \texttt{Table}, \texttt{Tree}, \texttt{Calendar}, and \texttt{Breadcrumbs}.
\item[\textbf{cssClass}]: An optional CSS class (the \texttt{class} attribute) to apply to the content element for styling.
\item[\textbf{section}]: The name of the page section where this content should be rendered (e.g., "header", "main", "sidebar").
\item[\textbf{sequence}]: A number that determines the order of content within a section.
\item[\textbf{dsn}]: The \texttt{datasource} attribute on the \texttt{query} element, which identifies a specific data source configured in the webbase root.
\item[\textbf{command}]: The text of the \texttt{query} element, executed against the specified data source. Before execution, this text is processed by the \textbf{Webbase Placeholders Language (\wbpl{})} processor. This processor replaces placeholders within the command with values sourced from several locations: the \texttt{parameters} element, the URL querystring, session variables, global variables, and HTTP headers. After this substitution, the resulting command is executed by the data source using its native language (e.g., SQL for a relational database, or JSONata if the data source is a JSON-based API).
\item[\textbf{params}]: The \texttt{parameters} element, a string formatted as a standard querystring (e.g., \texttt{key1=value1\&key2=value2}).
\item[\textbf{layout}]: The \texttt{layouts} element that defines, per language, how the fetched data should be rendered.
\end{description}

\subsubsection{Content Categories}
\label{sec:content-categories}

Contents are the fundamental elements of \wbol{} and represent interactive data units that fall into four distinct categories, each serving a specific purpose in the portal's user interface:

\begin{description}
\item[\textbf{Sensorial Contents}]: These render data in human-readable formats for information consumption and input. They include:
\begin{itemize}
\item Free text displays
\item Forms for data input
\item Lists and tables for structured data presentation
\item Plots and charts for data visualization
\item Maps for geographical data
\item Timelines for temporal data
\item Media elements (images, videos, audio)
\end{itemize}

\item[\textbf{Navigational Contents}]: These render data as interactive links and navigation elements, facilitating movement through the portal structure:
\begin{itemize}
\item Menus (primary and secondary navigation)
\item Table of contents (TOC)
\item Breadcrumbs for hierarchical navigation
\item Slicers for data filtering
\item Image maps with clickable regions
\item Pagination controls
\end{itemize}

\item[\textbf{Organizational Contents}]: These wrap and organize other contents in structured manners, providing logical grouping and presentation:
\begin{itemize}
\item Tabs for content organization
\item Calendars for date-based content
\item Trees for hierarchical data
\item Graphs for relationship visualization
\item Accordions for collapsible sections
\item Carousels for sequential content display
\end{itemize}

\item[\textbf{Special Contents}]: These provide specialized functionalities:
\begin{itemize}
\item Shortcuts are contents that point to other contents. It's a form of content reuse.
\item Client-side code are execution blocks (JavaScript)
\item Server-side code are execution blocks (API)
\end{itemize}
\end{description}

All content types explicitly declare or request the data they interact with, and their rendering behavior is determined by their category and specific subtype. The content categorization ensures consistency in user interface patterns and enables the Web Spinner to apply appropriate rendering logic and security policies.

\subsection{STWContentWithOptions}

\texttt{STWContentWithOptions} is similar to \texttt{STWContent}, it contains a list of \texttt{option} elements (GUIDs). This type is useful for scenarios like menus and tabs where the content is sourced from other elements.

\subsubsection{Property Description}

\begin{description}
\item[\textbf{subtype}]: Specifies the type of content to be rendered. Possible values include \texttt{Menu}, \texttt{Tabs}, and \texttt{Accordion}.
\end{description}

\section{Looking Forward}
\label{sec:wbol-forward}

A static portal, however well-structured, has limited utility in a dynamic business environment. The true power of an enterprise portal lies in its ability to present real-time, contextualized data. But how do we bridge the gap between the static structure of \wbol{} and the dynamic data living in enterprise systems?

The next chapter introduces the \textbf{Webbase Placeholders Language (\wbpl{})}, a specialized language designed to embed dynamic queries and data-driven logic directly within your \wbol{} definitions. We will explore how \wbpl{} turns your static templates into living, breathing documents.
